En SOAP todo va en un sobre XML con una cabecera opcional y un cuerpo; si algo falla, el cuerpo trae un \texttt{Fault} \cite{w3c_soap12}. El contrato es el WSDL, que fija operaciones, mensajes y dirección \cite{w3c_wsdl11}. El SRI es un buen ejemplo: su operación \texttt{validarComprobante} recibe la factura en XML codificado en base64 \cite{sri_recepcion}.

En REST cada recurso tiene una URI y los verbos HTTP dicen qué se hace con él: \texttt{GET} lee, \texttt{POST} crea, \texttt{PUT} reemplaza y \texttt{DELETE} borra \cite{rfc9110}. El servidor no guarda estado entre peticiones \cite{fielding2000} y el contrato, si existe, se escribe en OpenAPI \cite{openapi31}. En el diagrama se ve nuestro PFC: Angular consulta el estado por REST con JSON y un sistema externo puede pedir lo mismo por SOAP; los dos caminos terminan en el mismo servicio de negocio y la misma base de datos.
